% GENERATED FILE. Edit canonical lessons, then run npm run generate:books.
% canonical-source-hash: fnv1a64:ad1b8c75c47c4541
% canonical-lessons: JA-C166-gozen, JA-C166-hiru, JA-C166-ban, JA-C166-konban, JA-C166-maiban

\chapter{Mornings, Noons and Nights Ahead}
\label{ch:ja-mornings-noons-and-nights-ahead}
\hlchaptermodality{166}

\begin{chapteropening}
\textbf{By the end of this chapter:} \emph{I can say what time of day I will do things.}

Complete the last lesson of chapter 166: I can say what time of day I will do things.
\end{chapteropening}

\section[\texorpdfstring{\ja{ごぜん} (gozen)}{ごぜん (gozen)}]{\ja{ごぜん} (gozen) — a.m.; the morning}

\label{lesson:JA-C166-gozen}



\begin{quote}
Before the new one: say the Japanese for the twentieth of the month, then the Japanese for a birthday.
\end{quote}

\subsection*{You'll want to know: \ja{ごぜん}}
\textbf{\ja{ごぜん}} — \emph{gozen} — \textquotedblleft{}a.m.; the morning\textquotedblright{}.

\begin{quote}
\textbf{\ja{あしたの} \ja{ごぜんは} \ja{はたらきます。}} — \emph{ashita no gozen wa hatarakimasu} — \textquotedblleft{}Tomorrow morning I will work.\textquotedblright{}
\end{quote}

\begin{cousinweb}
\textbf{\ja{ごぜん}} (\emph{gozen}) is Sino-Japanese (kango): two Chinese characters read with their on'yomi, readings borrowed from Chinese. They are \emph{go}, \textquotedblleft{}noon\textquotedblright{} (Middle Chinese \emph{nguX}) and \emph{zen}, \textquotedblleft{}before, in front\textquotedblright{} (\emph{dzen}) (KANJIDIC2; Middle Chinese in Baxter's transcription).
\end{cousinweb}

\begin{grammarlens}[title={What you've built}]
One more word for everyday life at A2.
\end{grammarlens}

\subsection*{Your turn}
\begin{itemize}
\raggedright
  \item \emph{Say it:} \emph{gozen}
  \item \emph{Say it:} \emph{gozen}, once more
  \item \emph{Say it:} \emph{ashita no gozen wa hatarakimasu} — Tomorrow morning I will work.
  \item \emph{Recall:} say the Japanese for the twentieth of the month, then the Japanese for a birthday, then say \emph{gozen} again
\end{itemize}

\begin{tcolorbox}[breakable,colback=teal!4,colframe=teal!35!black,title={Before you move on}]
What does \ja{ごぜん} mean? (\textquotedblleft{}A.m.; the morning\textquotedblright{}.) Say it once more.
\end{tcolorbox}
\section[\texorpdfstring{\ja{ひる} (hiru)}{ひる (hiru)}]{\ja{ひる} (hiru) — noon; the daytime}

\label{lesson:JA-C166-hiru}



\begin{quote}
Before the new one: say the Japanese for a birthday, then the Japanese for a.m..
\end{quote}

\subsection*{You'll want to know: \ja{ひる}}
\textbf{\ja{ひる}} — \emph{hiru} — \textquotedblleft{}noon; the daytime\textquotedblright{}.

\begin{quote}
\textbf{\ja{あしたは} \ja{ひるに} \ja{ともだちと} \ja{あいます。}} — \emph{ashita wa hiru ni tomodachi to aimasu} — \textquotedblleft{}Tomorrow I will meet a friend at noon.\textquotedblright{}
\end{quote}

\begin{cousinweb}
\textbf{\ja{ひる}} (\emph{hiru}) is a native Japanese word (wago), inherited rather than borrowed from Chinese; its kanji was attached later as a kun'yomi, and its older origin is not settled (JMdict, KANJIDIC2).
\end{cousinweb}

\begin{grammarlens}[title={What you've built}]
One more word for everyday life at A2.
\end{grammarlens}

\subsection*{Your turn}
\begin{itemize}
\raggedright
  \item \emph{Say it:} \emph{hiru}
  \item \emph{Say it:} \emph{hiru}, once more
  \item \emph{Say it:} \emph{ashita wa hiru ni tomodachi to aimasu} — Tomorrow I will meet a friend at noon.
  \item \emph{Recall:} say the Japanese for a birthday, then the Japanese for a.m., then say \emph{hiru} again
\end{itemize}

\begin{tcolorbox}[breakable,colback=teal!4,colframe=teal!35!black,title={Before you move on}]
What does \ja{ひる} mean? (\textquotedblleft{}Noon; the daytime\textquotedblright{}.) Say it once more.
\end{tcolorbox}
\section[\texorpdfstring{\ja{ばん} (ban)}{ばん (ban)}]{\ja{ばん} (ban) — night-time; the evening}

\label{lesson:JA-C166-ban}



\begin{quote}
Before the new one: say the Japanese for a.m., then the Japanese for noon.
\end{quote}

\subsection*{You'll want to know: \ja{ばん}}
\textbf{\ja{ばん}} — \emph{ban} — \textquotedblleft{}night-time; the evening\textquotedblright{}.

\begin{quote}
\textbf{\ja{あしたの} \ja{ばんは} \ja{うちで} \ja{ゆうはんを} \ja{たべます。}} — \emph{ashita no ban wa uchi de yūhan o tabemasu} — \textquotedblleft{}Tomorrow evening I will have dinner at home.\textquotedblright{}
\end{quote}

\begin{cousinweb}
\textbf{\ja{ばん}} (\emph{ban}) is Sino-Japanese (kango): one Chinese character read with its on'yomi, a reading borrowed from Chinese. It is \emph{ban}, \textquotedblleft{}nightfall, evening\textquotedblright{} (Middle Chinese \emph{mjonX}) (KANJIDIC2; Middle Chinese in Baxter's transcription).
\end{cousinweb}

\begin{grammarlens}[title={What you've built}]
One more word for everyday life at A2.
\end{grammarlens}

\subsection*{Your turn}
\begin{itemize}
\raggedright
  \item \emph{Say it:} \emph{ban}
  \item \emph{Say it:} \emph{ban}, once more
  \item \emph{Say it:} \emph{ashita no ban wa uchi de yūhan o tabemasu} — Tomorrow evening I will have dinner at home.
  \item \emph{Recall:} say the Japanese for a.m., then the Japanese for noon, then say \emph{ban} again
\end{itemize}

\begin{tcolorbox}[breakable,colback=teal!4,colframe=teal!35!black,title={Before you move on}]
What does \ja{ばん} mean? (\textquotedblleft{}Night-time; the evening\textquotedblright{}.) Say it once more.
\end{tcolorbox}
\section[\texorpdfstring{\ja{こんばん} (konban)}{こんばん (konban)}]{\ja{こんばん} (konban) — this evening}

\label{lesson:JA-C166-konban}



\begin{quote}
Before the new one: say the Japanese for noon, then the Japanese for night-time.
\end{quote}

\subsection*{You'll want to know: \ja{こんばん}}
\textbf{\ja{こんばん}} — \emph{konban} — \textquotedblleft{}this evening\textquotedblright{}.

\begin{quote}
\textbf{\ja{こんばん} \ja{ともだちと} \ja{あいます。}} — \emph{konban tomodachi to aimasu} — \textquotedblleft{}I am meeting a friend this evening.\textquotedblright{}
\end{quote}

\begin{cousinweb}
\textbf{\ja{こんばん}} (\emph{konban}) is Sino-Japanese (kango): two Chinese characters read with their on'yomi, readings borrowed from Chinese. They are \emph{kon}, \textquotedblleft{}now\textquotedblright{} (Middle Chinese \emph{kim}) and \emph{ban}, \textquotedblleft{}nightfall, evening\textquotedblright{} (\emph{mjonX}) (KANJIDIC2; Middle Chinese in Baxter's transcription).
\end{cousinweb}

\begin{grammarlens}[title={What you've built}]
One more word for everyday life at A2.
\end{grammarlens}

\subsection*{Your turn}
\begin{itemize}
\raggedright
  \item \emph{Say it:} \emph{konban}
  \item \emph{Say it:} \emph{konban}, once more
  \item \emph{Say it:} \emph{konban tomodachi to aimasu} — I am meeting a friend this evening.
  \item \emph{Recall:} say the Japanese for noon, then the Japanese for night-time, then say \emph{konban} again
\end{itemize}

\begin{tcolorbox}[breakable,colback=teal!4,colframe=teal!35!black,title={Before you move on}]
What does \ja{こんばん} mean? (\textquotedblleft{}This evening\textquotedblright{}.) Say it once more.
\end{tcolorbox}
\section[\texorpdfstring{\ja{まいばん} (maiban)}{まいばん (maiban)}]{\ja{まいばん} (maiban) — every night}

\label{lesson:JA-C166-maiban}



\begin{quote}
Before the new one: say the Japanese for night-time, then the Japanese for this evening.
\end{quote}

\subsection*{You'll want to know: \ja{まいばん}}
\textbf{\ja{まいばん}} — \emph{maiban} — \textquotedblleft{}every night\textquotedblright{}.

\begin{quote}
\textbf{\ja{らいしゅうは} \ja{まいばん} \ja{はたらきます。}} — \emph{raishū wa maiban hatarakimasu} — \textquotedblleft{}Next week I will work every night.\textquotedblright{}
\end{quote}

\begin{cousinweb}
\textbf{\ja{まいばん}} (\emph{maiban}) is Sino-Japanese (kango): two Chinese characters read with their on'yomi, readings borrowed from Chinese. They are \emph{mai}, \textquotedblleft{}every\textquotedblright{} (Middle Chinese \emph{mwojX}) and \emph{ban}, \textquotedblleft{}nightfall, evening\textquotedblright{} (\emph{mjonX}) (KANJIDIC2; Middle Chinese in Baxter's transcription).
\end{cousinweb}

\begin{grammarlens}[title={What you've built}]
One more word for everyday life at A2.
\end{grammarlens}

\subsection*{Your turn}
\begin{itemize}
\raggedright
  \item \emph{Say it:} \emph{maiban}
  \item \emph{Say it:} \emph{maiban}, once more
  \item \emph{Say it:} \emph{raishū wa maiban hatarakimasu} — Next week I will work every night.
  \item \emph{Recall:} say the Japanese for night-time, then the Japanese for this evening, then say \emph{maiban} again
\end{itemize}

\begin{tcolorbox}[breakable,colback=teal!4,colframe=teal!35!black,title={Before you move on}]
What does \ja{まいばん} mean? (\textquotedblleft{}Every night\textquotedblright{}.) Say it once more.
\end{tcolorbox}
